% !TeX root = ../main.tex
\section{Device State and Working Set}\label{sec:device-state}

\begin{hypothesis}
  Back-to-back sustained sequential writes should first ride the SLC/DRAM write cache (a burst), then
  settle to a lower steady-state throughput once the cache must flush to QLC in the background. Occasional
  latency spikes from garbage collection should appear in the upper tail even while mean throughput is
  steady.
\end{hypothesis}

\begin{hypothesis}
  (Added before the large-file runs.) (i) A single sequential write far larger than the SLC cache should
  show a sharp \emph{cliff}: burst throughput near the QD=32 write rate of \cref{sec:block-size}, then a
  drop to the native QLC program rate, several times lower. (ii) After that, random \qty{4}{\kibi\byte}
  writes reach a lower steady state than in the burst, with heavier tails from garbage collection.
  (iii) A short idle period lets the controller drain the cache and recover part of the burst rate.
  (iv) Random reads should slow as the working set grows from \qty{1}{\gibi\byte} to
  \qty{60}{\gibi\byte}, because more logical-to-physical map lookups miss the controller's mapping cache.
  (v) \enquote{Sequential} reads of data written randomly should be much slower than of data written
  sequentially. (vi) Temperature should rise during sustained writes but stay below the
  \qty{86.85}{\celsius} warning threshold reported in sysfs (\code{temp1_max}).
\end{hypothesis}

\begin{knobs}
  \code{sustained} & sequential \qty{128}{\kibi\byte} QD32 fill of a fresh \qty{60}{\gibi\byte} file
                     (time cap \qty{3600}{\second}); then \qty{4}{\kibi\byte} QD32 random writes over the
                     whole file (\qty{180}{\second} or \qty{8}{\gibi\byte}), a \qty{60}{\second} idle gap and a
                     \qty{128}{\kibi\byte} QD32 recovery probe (\qty{15}{\second} or \qty{4}{\gibi\byte});
                     logged every second \\
  \code{working_set} & after a \qty{20}{\minute} idle period, on the same \qty{60}{\gibi\byte} file:
                     \qty{4}{\kibi\byte} random read (QD1, QD32) and write (QD32, \qty{256}{\mebi\byte} cap)
                     over its first 1, 4, 16 and \qty{60}{\gibi\byte}, then in reverse order, 3 repetitions
                     per pass \\
  \code{layout}     & first run (\code{20261021-103201}), \qty{2}{\gibi\byte} file: \qty{128}{\kibi\byte}
                     sequential (QD1, QD32) and \qty{4}{\kibi\byte} random reads after a sequential fill,
                     then after a \qty{4}{\gibi\byte} random overwrite (\ensuremath{\sim}\qty{86}{\percent} of
                     blocks) \\
  Drive state       & \code{fstrim} of \qty{42.8}{\gibi\byte} at about 01:41, then \qty{30}{\minute} idle;
                     fill at 02:11. \qty{77}{\gibi\byte} free on the scratch partition; other partitions
                     (one untouched \qty{351}{\gibi\byte}) not trimmed \\
  Script            & \code{rerun_clean.sh} (\code{sustained}, \code{working_set}); \code{run.sh layout} \\
\end{knobs}

\subsection{SLC-cache exhaustion and steady state}
\pgfplotsset{susplot/.style={acs, height=0.62\linewidth, legend style={font=\scriptsize}}}
\begin{figure}[!htbp]
  \centering
  \IfDataExists{sustained}{%
    \begin{subfigure}[t]{0.49\linewidth}
      \begin{tikzpicture}
        \begin{axis}[susplot, xlabel={GiB written}, ylabel={Write MiB/s (1 s bins)}, ymin=0]
          \addplot[mark=none, blue!70!black] table[x=gib_done, y=mibps, only if={phase}{seq_fill}] {data/sustained.csv};
        \end{axis}
      \end{tikzpicture}
      \caption{Sequential fill: throughput vs.\ volume.}
    \end{subfigure}\hfill
    \begin{subfigure}[t]{0.49\linewidth}
      \begin{tikzpicture}
        \begin{axis}[susplot, xlabel={GiB written}, ylabel={Latency (\si{\micro\second})}, ymode=log,
                     legend pos=north west]
          \addplot[mark=none, black!60] table[x=gib_done, y=lat_p50_us, only if={phase}{seq_fill}] {data/sustained.csv};
          \addlegendentry{p50}
          \addplot[mark=none, red!70!black] table[x=gib_done, y=lat_p999_us, only if={phase}{seq_fill}] {data/sustained.csv};
          \addlegendentry{p99.9}
        \end{axis}
      \end{tikzpicture}
      \caption{Sequential fill: per-second latency.}
    \end{subfigure}

    \begin{subfigure}[t]{0.49\linewidth}
      \begin{tikzpicture}
        \begin{axis}[susplot, xlabel={Time in phase (s)}, ylabel={4 KiB write IOPS}, ymin=0]
          \addplot[mark=none, blue!70!black] table[x=t_s, y=iops, only if={phase}{rand_steady}] {data/sustained.csv};
        \end{axis}
      \end{tikzpicture}
      \caption{Random-write phase (QD32).}
    \end{subfigure}\hfill
    \begin{subfigure}[t]{0.49\linewidth}
      \begin{tikzpicture}
        \begin{axis}[susplot, xlabel={Time (s)}, ylabel={NVMe temperature (\si{\celsius})},
                     restrict y to domain=0:120]
          \addplot[mark=none, orange!80!black] table[x=t_global_s, y=temp_c] {data/sustained.csv};
        \end{axis}
      \end{tikzpicture}
      \caption{Composite temperature, all phases.}
    \end{subfigure}%
  }
  \caption{Sustained-write experiment, one-second samples. (a) The SLC-cache cliff appears as the drop in
  throughput after \csvval[1]{sustained_summary}{cliff_gib}\,GiB; the fill then wrote the rest of the
  \csvval[1]{sustained_summary}{fill_gib}\,GiB at the QLC rate. (b) The latency during the fill. (c) The
  random-write phase after the fill, which reached its \qty{8}{\gibi\byte} cap after \qty{153}{\second}.
  (d) Temperature across the baseline idle, fill, random, idle and recovery phases.}\label{fig:sustained}
\end{figure}

\begin{table}[!htbp]
  \centering\small
  \caption{Sustained-write summary (computed by \code{run.sh} into \code{sustained_summary.csv}). Burst is
  the median of the leading seconds that stay above half the first second's rate. The cliff is the first
  second, starting a 5-second window, whose rate falls below half the burst; the table gives the volume
  written before it. Steady is the median of all seconds after the cliff. For the random phase the table
  gives the byte-weighted average, the per-second quartiles and the median per-second p99.}\label{tab:sustained}
  \begin{tabular}{@{}lr@{}}
    \toprule
    Burst sequential write (MiB/s) / duration (s)  & \csvval{sustained_summary}{burst_mibps} / \csvval{sustained_summary}{burst_s} \\
    Cliff after (GiB written / s)                  & \ssdcliffcell \\
    Post-cliff steady sequential write (MiB/s)     & \ssdifcliff{\csvval{sustained_summary}{steady_mibps}}{n/a} \\
    Steady / burst                                 & \ssdifcliff{\csvval[3]{sustained_summary}{steady_over_burst}}{n/a} \\
    Fill: GiB, seconds, average MiB/s              & \csvval[1]{sustained_summary}{fill_gib}, \csvval{sustained_summary}{fill_s}, \csvval{sustained_summary}{fill_avg_mibps} \\
    Random 4 KiB write, QD32: phase average IOPS   & \csvval{sustained_summary}{rand_steady_iops} \\
    \quad per-second IOPS quartiles; median per-second p99 (\si{\micro\second})
                                                   & \csvval{sustained_summary}{rand_iops_q25}--\csvval{sustained_summary}{rand_iops_q75}; \csvval{sustained_summary}{rand_steady_p99_us} \\
    Recovery after \qty{60}{\second} idle: first / last \qty{5}{\second} (MiB/s) & \csvval{sustained_summary}{recovery_mibps} / \csvval{sustained_summary}{recovery_late_mibps} \\
    NVMe temperature min--max (\si{\celsius})      & \csvval[1]{sustained_summary}{temp_min_c}--\csvval[1]{sustained_summary}{temp_max_c} \\
    \bottomrule
  \end{tabular}
\end{table}

\paragraph{Analysis.} The fill ran at the SLC ingest rate,
\csvval{sustained_summary}{burst_mibps}\,MiB/s, for \csvval{sustained_summary}{burst_s}\,s and collapsed
after \csvval[1]{sustained_summary}{cliff_gib}\,GiB, within two seconds, to
\qtyrange{60}{69}{\mebi\byte\per\second}; the median per-second latency jumped from
\qty{3.4}{} to \ensuremath{\sim}\qty{63}{\milli\second}. The post-cliff median,
\csvval{sustained_summary}{steady_mibps}\,MiB/s (\csvval[3]{sustained_summary}{steady_over_burst} of the
burst, a \ssdburstratio$\times$ drop, far more than the \enquote{several times} of hypothesis~(i)), held
for the remaining \qty{39}{\gibi\byte} with no downward drift: per-minute medians stayed at
\qtyrange{62}{64}{\mebi\byte\per\second} except one minute, near \qty{53}{\gibi\byte}, dominated by a
slow episode. The rate was not smooth, though: \qty{69}{\percent} of post-cliff seconds ran at
\qtyrange{60}{70}{\mebi\byte\per\second} and \qty{19}{\percent} below \qty{20}{\mebi\byte\per\second}.
The whole \qty{60}{\gibi\byte} took \csvval{sustained_summary}{fill_s}\,s
(\csvval{sustained_summary}{fill_avg_mibps}\,MiB/s on average). We take \ensuremath{\sim}\qty{63}{\mebi\byte\per\second}
as the direct-to-QLC rate while the drive folds; hypothesis~(i) holds. One detail we cannot explain: for
four seconds, between \qty{3.1}{} and \qty{5.7}{\gibi\byte} written, the burst dipped to \ensuremath{\sim}\qty{670}{\mebi\byte\per\second}
before full speed returned.

\textbf{How large is the cache? It depends on drive state.} After trimming \qty{42.8}{\gibi\byte} and
\qty{30}{\minute} idle the cliff came at \csvval[1]{sustained_summary}{cliff_gib}\,GiB, close to the
\ensuremath{\sim}\qty{20}{\gibi\byte} we expected. Two superseded runs found much less: \qty{3.8}{\gibi\byte}
in the lean run (\code{sustained/20261021-220614}, also after \code{fstrim} and $\geq$\qty{30}{\minute} idle,
but the same day as a \qty{60}{\gibi\byte} fill), and \qty{3.1}{\gibi\byte} in the untrimmed first run
(\code{sustained/20261021-105439}). A dynamic SLC cache sized by the free blocks the controller has
available, common in QLC client drives \cite{tallis2018intel660p,micron2210brief}, would explain all
three, but we did not measure the free-block count, so this remains an untested hypothesis.

The random-write phase (\qty{4}{\kibi\byte}, QD32, over the whole file) reached its
\qty{8}{\gibi\byte} cap after \qty{153}{\second} at \csvval{sustained_summary}{rand_steady_iops} IOPS on
average (per-second quartiles \csvval{sustained_summary}{rand_iops_q25}--\csvval{sustained_summary}{rand_iops_q75},
median per-second p99 \qty{33.8}{\milli\second}), about $18\times$ below the in-cache rate
(\csvq{queue_depth}{qd=32,mode=write}{iops_median} IOPS). Most seconds ran at
\ensuremath{\sim}\num{16}k IOPS (\ensuremath{\sim}\qty{63}{\mebi\byte\per\second}, the same QLC rate as
the sequential fill), and in three episodes (34 of 154 seconds) at \ensuremath{\sim}\num{3}k IOPS. This
supports hypothesis~(ii): a much lower rate with heavy tails, sustained for \qty{2.5}{\minute}.

After the \qty{60}{\second} idle gap the recovery probe wrote at
\csvval{sustained_summary}{recovery_mibps}\,MiB/s (last \qty{5}{\second}:
\csvval{sustained_summary}{recovery_late_mibps}\,MiB/s), the QLC rate: no recovery, as in the first run
(\qty{61}{} then \qty{19}{\mebi\byte\per\second}). Hypothesis~(iii) fails after a full
\qty{60}{\gibi\byte} fill. In the lean run, after a fill of only \qty{4.3}{\gibi\byte}, the same gap
restored the burst rate for at least \qty{1}{\gibi\byte} (the probe's cap), so recovery depends on the
fold backlog. Even the \qty{20}{\minute} idle period before \code{working_set} freed only about
\qty{1.2}{\gibi\byte} of fast writes (\cref{sec:ws}). This is short-term behaviour, not the full-drive SNIA
steady state \cite{snia_pts}.

\paragraph{Temperature.} The composite temperature stayed at
\csvval[1]{sustained_summary}{temp_min_c}--\csvval[1]{sustained_summary}{temp_max_c}\,\si{\celsius} here and
\qtyrange{31.9}{60.9}{\celsius} over the whole run, far below the \qty{86.85}{\celsius} warning threshold,
so hypothesis~(vi) holds. The burst ran at up to \qty{42.9}{\celsius} before the cliff, and the recovery
probe still ran at the QLC rate at \qtyrange{41.9}{44.9}{\celsius}, so the cliff is the cache, not thermal
throttling.

\subsection{Working set (LBA range)}\label{sec:ws}
\begin{figure}[!htbp]
  \centering
  \IfDataExists{working_set}{%
    \begin{subfigure}[t]{0.49\linewidth}
      \begin{tikzpicture}
        \begin{axis}[acs, height=0.85\linewidth, xmode=log, log basis x=2, xtick=data, log ticks with fixed point,
                     xlabel={Working set (GiB)}, ylabel={QD=1 read mean latency (\si{\micro\second})}]
          \addplot+[iqr bars] table[x=span_gib, y=lat_mean_us_median, iqr of=lat_mean_us, only if={mode}{read}, only if={qd}{1}] {data/working_set.csv};
        \end{axis}
      \end{tikzpicture}
      \caption{QD=1 random read latency.}
    \end{subfigure}\hfill
    \begin{subfigure}[t]{0.49\linewidth}
      \begin{tikzpicture}
        \begin{axis}[acs, height=0.85\linewidth, xmode=log, log basis x=2, xtick=data, log ticks with fixed point,
                     xlabel={Working set (GiB)}, ylabel={QD=32 IOPS}, ymin=0, legend pos=south west]
          \addplot+[iqr bars] table[x=span_gib, y=iops_median, iqr of=iops, only if={mode}{read}, only if={qd}{32}] {data/working_set.csv};
          \addlegendentry{rand read}
          \addplot+[dashed, iqr bars] table[x=span_gib, y=iops_median, iqr of=iops, only if={mode}{write}, only if={qd}{32}] {data/working_set.csv};
          \addlegendentry{rand write}
        \end{axis}
      \end{tikzpicture}
      \caption{QD=32 random read and write IOPS.}
    \end{subfigure}%
  }
  \caption{\qty{4}{\kibi\byte} random I/O restricted to the first 1, 4, 16 and \qty{60}{\gibi\byte} of the
  \qty{60}{\gibi\byte} file after a \qty{20}{\minute} idle period (medians of 6 repetitions, 3 in each of
  a forward and a reverse pass, with IQR bars). QD=32 read IOPS halve between 1 and
  \qty{4}{\gibi\byte} and then stay flat; beyond \qty{4}{\gibi\byte} QD=1 latency differs between the
  passes (wide IQRs), and the write IQRs span a cache-exhaustion episode (see text).}\label{fig:working-set}
\end{figure}

\textbf{Reads.} QD=32 read IOPS halve between 1 and \qty{4}{\gibi\byte}
(\csvq{working_set}{span_gib=1,mode=read,qd=32}{iops_median}$\to$%
\csvq{working_set}{span_gib=4,mode=read,qd=32}{iops_median}) in both passes and then stay flat
(\csvq{working_set}{span_gib=16,mode=read,qd=32}{iops_median} and
\csvq{working_set}{span_gib=60,mode=read,qd=32}{iops_median} at 16 and \qty{60}{\gibi\byte}). QD=1
latency also rises from 1 to \qty{4}{\gibi\byte}
(\csvq[1]{working_set}{span_gib=1,mode=read,qd=1}{lat_mean_us_median}$\to$%
\csvq[1]{working_set}{span_gib=4,mode=read,qd=1}{lat_mean_us_median}\,\si{\micro\second} mean), but
beyond that it depends on the pass more than on the span: \qty{96}{} and \qty{106}{\micro\second} forward at 16
and \qty{60}{\gibi\byte}, \qty{55.5}{\micro\second} at both in the reverse pass, hence the wide IQRs.

Hypothesis~(iv) predicted a steady rise as more map lookups miss the controller's mapping cache. The data
do not show it: nothing grows from 4 to \qty{60}{\gibi\byte}, although the map for the span grows from
\qty{4}{} to \qty{60}{\mebi\byte} (4 bytes per \qty{4}{\kibi\byte} page). The step between 1 and
\qty{4}{\gibi\byte} did not appear in the lean run's 1--\qty{4}{\gibi\byte} spans on a freshly written
\qty{4}{\gibi\byte} file (flat QD=1 latency, QD=32 IOPS within \qty{0.3}{\percent}), so it reflects this
file's write history (a sequential fill mostly past the SLC cache, then \qty{8}{\gibi\byte} of random
overwrites) more than map size. Even at 4--\qty{60}{\gibi\byte} the file gives only about
\qty{40}{\percent} of the QD=32 read rate of the freshly written \qty{2}{\gibi\byte} test file
(\csvq{queue_depth}{qd=32,mode=read}{iops_median} IOPS), again pointing to where the data sit; why the
first gigabyte is faster is untested. We count hypothesis~(iv) as not supported.

\textbf{Writes track time, not span.} QD=32 writes ran at the in-cache rate
(\ensuremath{\sim}\qty{950}{\mebi\byte\per\second}) for the first \ensuremath{\sim}\qty{1.2}{\gibi\byte}
(the \qty{1}{\gibi\byte} span and the start of the \qty{4}{\gibi\byte} one), fell to
\qtyrange{61}{67}{\mebi\byte\per\second} at 16 and \qty{60}{\gibi\byte} (as low as
\qty{19}{\mebi\byte\per\second} at the start of the reverse pass), and returned to
\ensuremath{\sim}\qty{950}{\mebi\byte\per\second} from the second reverse \qty{16}{\gibi\byte}
repetition on. The medians in \cref{fig:working-set}b therefore show the cache state at each moment;
the 20-minute idle period had freed only a little SLC space after the fill.

\subsection{Data layout: why preconditioning order matters}\label{sec:layout}
\begin{table}[!htbp]
  \centering\small
  \caption{The same reads on the same \qty{2}{\gibi\byte} file after a sequential fill and after a random
  \qty{4}{\kibi\byte} overwrite (medians of 5 repetitions; first run, \code{20261021-103201}).}\label{tab:layout}
  \begin{tabular}{@{}lrr@{}}
    \toprule
    Read & after sequential fill & after random overwrite \\
    \midrule
    \qty{128}{\kibi\byte} seq, QD=1 (MiB/s)  & \csvq{layout}{experiment=layout_seq_filled,bs=131072,qd=1}{mibps_median}
                                              & \csvq{layout}{experiment=layout_rand_overwritten,bs=131072,qd=1}{mibps_median} \\
    \qty{128}{\kibi\byte} seq, QD=32 (MiB/s) & \csvq{layout}{experiment=layout_seq_filled,bs=131072,qd=32}{mibps_median}
                                              & \csvq{layout}{experiment=layout_rand_overwritten,bs=131072,qd=32}{mibps_median} \\
    \qty{4}{\kibi\byte} rand, QD=1 (IOPS)    & \csvq{layout}{experiment=layout_seq_filled,bs=4096}{iops_median}
                                              & \csvq{layout}{experiment=layout_rand_overwritten,bs=4096}{iops_median} \\
    \bottomrule
  \end{tabular}
\end{table}

An SSD places data in the order it was \emph{written}, not by logical address, and \cref{tab:layout}
confirms hypothesis~(v) strongly. After the random overwrite, QD=1 sequential reads drop from
\csvq{layout}{experiment=layout_seq_filled,bs=131072,qd=1}{mibps_median} to
\csvq{layout}{experiment=layout_rand_overwritten,bs=131072,qd=1}{mibps_median}\,MiB/s ($6.4\times$), while
\qty{4}{\kibi\byte} random reads are unaffected. Our sub-prediction that QD=32 would hide the gap was
\emph{wrong}: it grows to $9.2\times$. Each fragmented \qty{128}{\kibi\byte} request becomes up to 32
scattered page reads, and the drive completes \ensuremath{\sim}\num{1830} such requests per second (about
\num{58}k page reads/s) at either depth, well below the \ensuremath{\sim}\num{136}k random reads/s it
sustains at QD=32 (\cref{sec:queue-depth}); why is untested. A sequential benchmark is meaningful only if the data's write history is controlled, which
is why \code{run.sh} rewrites the file before every sequential-read group.
